% =============================================================================
% CHAPTER 4 — IMPLEMENTATION
% =============================================================================
\chapter{Implementation}
\label{ch:impl}

\section{Technology Stack and Justification}

The technology choices for StockSense AI were governed by three criteria: cost, accessibility,
and capability. Each component of the stack was selected because it satisfies all three criteria
simultaneously. Table~\ref{tab:stack} presents the complete technology stack.

{\singlespacing
\renewcommand{\arraystretch}{1.5}
\begin{longtable}{@{} >{\raggedright\arraybackslash}p{3.5cm} >{\raggedright\arraybackslash}p{2.5cm} >{\raggedright\arraybackslash}p{8.5cm} @{}}
\caption{Technology Stack and Justification}
\label{tab:stack} \\
\toprule
\textbf{Technology} & \textbf{Role} & \textbf{Justification} \\
\midrule
\endfirsthead
\toprule
\textbf{Technology} & \textbf{Role} & \textbf{Justification} \\
\midrule
\endhead
\bottomrule
\endfoot
React 19 & Frontend framework &
  Component-based architecture; Concurrent Mode rendering for real-time UI updates;
  large ecosystem and strong community support. \\
Vite 8 & Build tool &
  Sub-second hot module replacement; built-in reverse proxy eliminates CORS issues
  in development; fast production bundle generation. \\
Tailwind CSS v4 & Styling &
  Utility-first CSS with the new \texttt{@theme} directive eliminates need for a
  separate configuration file; consistent design tokens across all components. \\
React Router 7 & Client routing &
  Declarative route definitions; nested layouts with \texttt{Outlet}; programmatic
  navigation for post-auth redirection. \\
Supabase & Database + Auth &
  Managed PostgreSQL with built-in authentication, real-time WebSocket subscriptions,
  and a JavaScript client; generous free tier; no server management required. \\
n8n (self-hosted) & Workflow engine &
  Low-code automation with visual workflow editor; native HTTP, Supabase, Groq,
  and Twilio integrations; cron scheduler for daily orchestration. \\
Groq Llama 3.1 8B & LLM inference &
  Fastest available inference speed for the Llama~3.1 family; free tier with
  generous rate limits suitable for prototype; JSON output mode for reliable parsing. \\
Google Gemini 2.5 Flash & Multimodal AI &
  State-of-the-art vision-language model for image understanding; supports base64
  image input; reliable JSON structured output for product and invoice extraction. \\
Twilio WhatsApp API & Alert delivery &
  Industry-standard CPaaS with sandbox WhatsApp support; simple REST API; per-message
  pricing model viable at kirana store scale. \\
ngrok & Tunnelling &
  HTTPS tunnel for exposing local n8n instance; used in development and
  early production; replaced by cloud hosting for persistent deployment. \\
\end{longtable}
}

\section{Development Environment}

The development environment was configured as follows:

\begin{itemize}[itemsep=3pt]
  \item \textbf{Operating System:} Windows 11 with PowerShell 7
  \item \textbf{Node.js:} v22 LTS (required for Vite 8 and React 19)
  \item \textbf{Package Manager:} npm with \texttt{--legacy-peer-deps} flag due to
        React 19 peer dependency resolution requirements
  \item \textbf{n8n:} Self-hosted via \texttt{npx n8n} on \texttt{localhost:5678}
  \item \textbf{Supabase:} Cloud-hosted project (free tier) accessed via the JavaScript
        client library \texttt{@supabase/supabase-js}
  \item \textbf{Environment Variables:} Stored in \texttt{frontend/.env} with
        \texttt{VITE\_} prefix for Vite client-side injection:
        \texttt{VITE\_SUPABASE\_URL}, \texttt{VITE\_SUPABASE\_ANON\_KEY},
        \texttt{VITE\_N8N\_BASE\_URL}, \texttt{VITE\_N8N\_API\_KEY}
  \item \textbf{Code Editor:} Visual Studio Code with ESLint extension
  \item \textbf{ESLint:} Flat config (ESLint 9) with
        \texttt{eslint-plugin-react-hooks} and \texttt{eslint-plugin-react-refresh}
  \item \textbf{Version Control:} Git with GitHub remote; single \texttt{main} branch;
        Conventional Commits message format
\end{itemize}

The Vite development server is configured to proxy all requests matching \texttt{/webhook/*}
and \texttt{/api/v1/*} to the n8n instance URL defined in \texttt{VITE\_N8N\_BASE\_URL},
eliminating browser-imposed CORS rejections without requiring changes to n8n's own CORS
configuration.

\section{Module-wise Implementation}

\subsection{WF-01: Product and Sales Data Ingestion}

WF-01 exposes a POST webhook at \texttt{/webhook/product-sales-ingest}. Upon receiving a
request from the frontend Add Product form, the workflow cleans and normalises the incoming
fields (removing whitespace, setting default values for optional fields), performs a lookup of
the Products table by \texttt{product\_id}, and branches on existence. If the product already
exists, it executes an UPDATE operation; if not, it inserts a new row. The webhook then responds
with a success payload that the frontend uses to dismiss the modal and trigger a real-time
refresh.

\begin{figure}[H]
\centering
\begin{tikzpicture}[node distance=0.5cm and 0.9cm, align=center]
  \node (n1) [wftrigger] {Webhook\\Trigger};
  \node (n2) [wfnode, right=of n1] {Edit Fields\\(Clean)};
  \node (n3) [wfnode, right=of n2] {Get Rows\\(Lookup)};
  \node (n4) [wfdecision, right=of n3] {Exists?};
  \node (n5a) [wfnode, above right=0.7cm and 1.0cm of n4] {Update\\Row};
  \node (n5b) [wfnode, below right=0.7cm and 1.0cm of n4] {Create\\Row};
  \node (n6) [wfnode, right=3.5cm of n4] {Respond\\Webhook};
  \draw[wfarrow] (n1)--(n2); \draw[wfarrow] (n2)--(n3); \draw[wfarrow] (n3)--(n4);
  \draw[wfarrow] (n4) |- node[wflabel,pos=0.25,above]{Yes} (n5a);
  \draw[wfarrow] (n4) |- node[wflabel,pos=0.25,below]{No}  (n5b);
  \draw[wfarrow] (n5a) -| (n6); \draw[wfarrow] (n5b) -| (n6);
\end{tikzpicture}
\caption{WF-01: Product and Sales Data Ingestion Workflow.}
\label{fig:wf01}
\end{figure}

\subsection{WF-02: Stockout Date Calculator}

WF-02 fetches all products from the Products table and applies the stockout prediction
formula. For each product, the workflow computes:

\begin{equation}
d_{\text{stockout}} = \left\lfloor \frac{S_{\text{current}}}{V_{\text{daily}}} \right\rfloor, \quad
\text{date}_{\text{stockout}} = \text{date}_{\text{today}} + d_{\text{stockout}}
\label{eq:stockout}
\end{equation}

Edge cases are handled explicitly: if $V_{\text{daily}} = 0$, then $d_{\text{stockout}}$ is set
to 999 (interpreted as "no foreseeable stockout"); if $S_{\text{current}} = 0$, then
$d_{\text{stockout}} = 0$ (immediate stockout). Both \texttt{days\_to\_stockout} and
\texttt{estimated\_stockout\_date} are written back to each product row.

\begin{figure}[H]
\centering
\begin{tikzpicture}[node distance=0.5cm and 0.9cm, align=center]
  \node (n1) [wftrigger] {Trigger};
  \node (n2) [wfnode, right=of n1] {Fetch All\\Products};
  \node (n3) [wfnode, right=of n2] {Calculate\\Stockout Dates};
  \node (n4) [wfnode, right=of n3] {Update\\Products};
  \draw[wfarrow] (n1)--(n2); \draw[wfarrow] (n2)--(n3); \draw[wfarrow] (n3)--(n4);
\end{tikzpicture}
\caption{WF-02: Stockout Date Calculator Workflow.}
\label{fig:wf02}
\end{figure}

\subsection{WF-03: Inventory Processing and Stock Alerts}

WF-03 identifies all products where \texttt{current\_stock} $\leq$ \texttt{reorder\_threshold}.
It first deletes all existing Stock Alert rows for the store (a clean-slate refresh strategy),
then inserts fresh alert records for each at-risk product with status \texttt{ACTIVE}. This
approach avoids duplicate alert accumulation and ensures the alert feed always reflects the
current inventory state.

\begin{figure}[H]
\centering
\begin{tikzpicture}[node distance=0.5cm and 0.8cm, align=center]
  \node (n1) [wftrigger] {Trigger};
  \node (n2) [wfnode, right=of n1] {Fetch\\Products};
  \node (n3) [wfnode, right=of n2] {Filter\\At-Risk};
  \node (n4) [wfnode, right=of n3] {Delete Old\\Alerts};
  \node (n5) [wfnode, right=of n4] {Restore\\Context};
  \node (n6) [wfnode, right=of n5] {Create New\\Alerts};
  \draw[wfarrow] (n1)--(n2); \draw[wfarrow] (n2)--(n3);
  \draw[wfarrow] (n3)--(n4); \draw[wfarrow] (n4)--(n5); \draw[wfarrow] (n5)--(n6);
\end{tikzpicture}
\caption{WF-03: Inventory Processing and Stock Alert Generation Workflow.}
\label{fig:wf03}
\end{figure}

\subsection{WF-04: Stock Risk Classification}

WF-04 classifies each product into one of four risk tiers using the following piecewise function
applied to \texttt{days\_to\_stockout}:

\begin{equation}
\text{risk} =
\begin{cases}
\text{HIGH}    & \text{if } d_{\text{stockout}} \leq 3 \\
\text{MEDIUM}  & \text{if } 3 < d_{\text{stockout}} \leq 7 \\
\text{LOW}     & \text{if } d_{\text{stockout}} > 7 \\
\text{UNKNOWN} & \text{if stockout data unavailable}
\end{cases}
\label{eq:risk}
\end{equation}

The classified risk level is written back to the \texttt{risk\_level} column of each product row
and is used by the frontend dashboard to colour-code product cards and by WF-05 to filter
at-risk products for AI reorder suggestion generation.

\begin{figure}[H]
\centering
\begin{tikzpicture}[node distance=0.5cm and 0.9cm, align=center]
  \node (n1) [wftrigger] {Trigger};
  \node (n2) [wfnode, right=of n1] {Fetch All\\Products};
  \node (n3) [wfnode, right=of n2] {Classify\\Risk Tiers};
  \node (n4) [wfnode, right=of n3] {Update\\Products};
  \draw[wfarrow] (n1)--(n2); \draw[wfarrow] (n2)--(n3); \draw[wfarrow] (n3)--(n4);
\end{tikzpicture}
\caption{WF-04: Stock Risk Classification Workflow.}
\label{fig:wf04}
\end{figure}

\subsection{WF-05: AI Reorder Intelligence}

WF-05 constitutes the most sophisticated workflow in the pipeline. It fetches at-risk products
(risk level HIGH or MEDIUM) and all suppliers in parallel, merges the two data streams, and
constructs a structured zero-shot prompt for the Groq Llama~3.1~8B model. The prompt
instructs the model to act as an inventory manager for an Indian kirana store and to return
a JSON array of reorder suggestions, one per at-risk product, with fields:
\texttt{product\_id}, \texttt{supplier\_id}, \texttt{suggested\_qty}
($= 30 \times V_{\text{daily}}$), and \texttt{reason} (one-sentence justification).

\begin{figure}[H]
\centering
\begin{tikzpicture}[node distance=0.5cm and 0.6cm, align=center]
  \node (n1) [wftrigger] {Trigger};
  \node (n2a) [wfnode, above right=0.5cm and 0.7cm of n1] {Fetch At-Risk\\Products};
  \node (n2b) [wfnode, below right=0.5cm and 0.7cm of n1] {Fetch All\\Suppliers};
  \node (n3) [wfnode, right=2.4cm of n1] {Merge};
  \node (n4) [wfnode, right=of n3] {Build AI\\Prompt};
  \node (n5) [wfnode, right=of n4] {Groq API\\Call};
  \node (n6) [wfnode, right=of n5] {Parse AI\\Response};
  \node (n7) [wfnode, below=0.8cm of n6] {Delete Old\\Suggestions};
  \node (n8) [wfnode, left=of n7] {Restore\\Context};
  \node (n9) [wfnode, left=of n8] {Insert New\\Suggestions};
  \draw[wfarrow] (n1) |- (n2a); \draw[wfarrow] (n1) |- (n2b);
  \draw[wfarrow] (n2a) -| (n3); \draw[wfarrow] (n2b) -| (n3);
  \draw[wfarrow] (n3)--(n4); \draw[wfarrow] (n4)--(n5);
  \draw[wfarrow] (n5)--(n6); \draw[wfarrow] (n6)--(n7);
  \draw[wfarrow] (n7)--(n8); \draw[wfarrow] (n8)--(n9);
\end{tikzpicture}
\caption{WF-05: AI Reorder Intelligence Workflow.}
\label{fig:wf05}
\end{figure}

\subsection{WF-06: Supplier Scoring}

WF-06 evaluates all suppliers using the Weighted Sum Model. The composite score $C_s$ is
computed as:

\begin{equation}
C_s = (R \times 0.40 + P \times 0.35 + D \times 0.25) \times 100
\label{eq:wsm}
\end{equation}

where $R$ is the reliability score (0--1), $P$ is the price score (0--1), and $D$ is the
delivery speed score (min-max normalised: fastest supplier $= 1.0$, slowest $= 0.0$). Letter
grades are assigned as: A ($C_s \geq 80$), B ($65 \leq C_s < 80$), C ($50 \leq C_s < 65$),
D ($C_s < 50$). Both score and grade are written to the Suppliers table.

\begin{figure}[H]
\centering
\begin{tikzpicture}[node distance=0.5cm and 0.9cm, align=center]
  \node (n1) [wftrigger] {Trigger};
  \node (n2) [wfnode, right=of n1] {Fetch All\\Suppliers};
  \node (n3) [wfnode, right=of n2] {Calculate\\WSM Scores};
  \node (n4) [wfnode, right=of n3] {Assign\\Grades};
  \node (n5) [wfnode, right=of n4] {Update\\Suppliers};
  \draw[wfarrow] (n1)--(n2); \draw[wfarrow] (n2)--(n3);
  \draw[wfarrow] (n3)--(n4); \draw[wfarrow] (n4)--(n5);
\end{tikzpicture}
\caption{WF-06: Supplier Scoring Workflow.}
\label{fig:wf06}
\end{figure}

\subsection{WF-07: WhatsApp Alert Sender}

WF-07 queries the Stock Alerts table for all records with \texttt{alert\_status = 'ACTIVE'}.
It applies a 24-hour anti-spam filter by checking \texttt{last\_alerted\_at}: only products
not alerted in the last 24 hours proceed to the message builder. A formatted WhatsApp message
is constructed listing all qualifying products with their names, current stock, and threshold
values. The message is dispatched to the registered store WhatsApp numbers via the Twilio
Messages API. After successful delivery, \texttt{last\_alerted\_at} is updated for each
product.

\begin{figure}[H]
\centering
\begin{tikzpicture}[node distance=0.5cm and 0.8cm, align=center]
  \node (n1) [wftrigger] {Trigger};
  \node (n2) [wfnode, right=of n1] {Fetch Active\\Alerts};
  \node (n3) [wfnode, right=of n2] {Filter\\(24hr cooldown)};
  \node (n4) [wfnode, right=of n3] {Build\\Message};
  \node (n5) [wfnode, right=of n4] {Send via\\Twilio};
  \node (n6) [wfnode, right=of n5] {Update\\Timestamps};
  \draw[wfarrow] (n1)--(n2); \draw[wfarrow] (n2)--(n3); \draw[wfarrow] (n3)--(n4);
  \draw[wfarrow] (n4)--(n5); \draw[wfarrow] (n5)--(n6);
\end{tikzpicture}
\caption{WF-07: WhatsApp Alert Sender Workflow.}
\label{fig:wf07}
\end{figure}

\subsection{WF-08: Master Daily Orchestrator}

WF-08 is triggered by an n8n Cron node set to 8:00 AM IST daily. It begins by simulating the
day's sales by subtracting each product's \texttt{avg\_daily\_sales} from its
\texttt{current\_stock} (with a floor of zero). It then sequentially executes WF-02 through
WF-07 as sub-workflow calls, waiting for each to complete before proceeding. After WF-07
completes, WF-08 computes the daily health score using the inventory health formula and inserts
a Daily Snapshot record. Finally, it writes a completion entry to System Logs.

\begin{figure}[H]
\centering
\begin{tikzpicture}[node distance=0.45cm and 0.6cm, align=center]
  \node (t)   [wftrigger] {Cron / Manual\\Trigger};
  \node (sim) [wfnode, right=of t]  {Simulate\\Daily Sales};
  \node (w2)  [wfnode, right=of sim] {WF-02\\Stockout};
  \node (w3)  [wfnode, right=of w2]  {WF-03\\Alerts};
  \node (w4)  [wfnode, right=of w3]  {WF-04\\Risk};
  \node (w5)  [wfnode, below=0.7cm of w4] {WF-05\\AI Reorder};
  \node (w6)  [wfnode, left=of w5]  {WF-06\\Scoring};
  \node (w7)  [wfnode, left=of w6]  {WF-07\\WhatsApp};
  \node (hs)  [wfnode, left=of w7]  {Compute\\Health Score};
  \node (sn)  [wfnode, left=of hs]  {Save\\Snapshot};
  \node (lg)  [wfnode, below=0.7cm of sn] {Log\\Completion};
  \draw[wfarrow] (t)--(sim); \draw[wfarrow] (sim)--(w2); \draw[wfarrow] (w2)--(w3);
  \draw[wfarrow] (w3)--(w4); \draw[wfarrow] (w4)--(w5); \draw[wfarrow] (w5)--(w6);
  \draw[wfarrow] (w6)--(w7); \draw[wfarrow] (w7)--(hs); \draw[wfarrow] (hs)--(sn);
  \draw[wfarrow] (sn)--(lg);
\end{tikzpicture}
\caption{WF-08: Master Daily Orchestrator Workflow.}
\label{fig:wf08}
\end{figure}

\subsection{WF-09: AI Vision Product Onboarding}

WF-09 exposes a POST webhook that accepts a base64-encoded image of a product. The image is
passed directly to the Gemini~2.5 Flash multimodal API with a structured prompt requesting
extraction of \texttt{product\_name} and \texttt{category} in JSON format. A recursive parser
handles Gemini's response envelope to reliably extract the JSON payload regardless of any
markdown fencing or preamble. The workflow then queries the Products table to determine the
highest existing product ID suffix and returns an auto-incremented collision-free ID along
with the extracted fields to the frontend form.

\begin{figure}[H]
\centering
\resizebox{\textwidth}{!}{%
\begin{tikzpicture}[node distance=0.5cm and 0.8cm, align=center]
  \node (n1) [wftrigger] {Webhook\\Trigger};
  \node (n2) [wfnode, right=of n1] {Gemini 2.5\\Flash};
  \node (n3) [wfnode, right=of n2] {Parse JSON\\Response};
  \node (n4) [wfnode, right=of n3] {Fetch Existing\\IDs};
  \node (n5) [wfnode, right=of n4] {Generate\\Unique ID};
  \node (n6) [wfnode, right=of n5] {Respond to\\Webhook};
  \draw[wfarrow] (n1)--(n2); \draw[wfarrow] (n2)--(n3); \draw[wfarrow] (n3)--(n4);
  \draw[wfarrow] (n4)--(n5); \draw[wfarrow] (n5)--(n6);
\end{tikzpicture}%
}
\caption{WF-09: AI Vision Onboarding Workflow.}
\label{fig:wf09}
\end{figure}

\subsection{WF-10: Smart Scan for Supplier Invoices}

WF-10 processes photographs of physical supplier delivery bills. The image is submitted to
Gemini~2.5 Flash with a prompt requesting extraction of all line items as a JSON array of
\texttt{\{product\_name, quantity\}} objects. The extracted items are then matched against
the store's existing Products table using a fuzzy string similarity algorithm implemented in
an n8n Code node. The algorithm computes character-level similarity scores between each
extracted product name and all inventory product names, selecting the best match above a
defined confidence threshold. The matched product IDs and computed new stock levels are
returned to the frontend for the store owner to review and confirm.

\begin{figure}[H]
\centering
\resizebox{\textwidth}{!}{%
\begin{tikzpicture}[node distance=0.5cm and 0.8cm, align=center]
  \node (n1) [wftrigger] {Webhook\\Trigger};
  \node (n2) [wfnode, right=of n1] {Gemini 2.5\\Flash};
  \node (n3) [wfnode, right=of n2] {Parse JSON\\Items};
  \node (n4) [wfnode, right=of n3] {Fetch Current\\Inventory};
  \node (n5) [wfnode, right=of n4] {Fuzzy Match\\\& Calculate};
  \node (n6) [wfnode, right=of n5] {Respond to\\Webhook};
  \draw[wfarrow] (n1)--(n2); \draw[wfarrow] (n2)--(n3); \draw[wfarrow] (n3)--(n4);
  \draw[wfarrow] (n4)--(n5); \draw[wfarrow] (n5)--(n6);
\end{tikzpicture}%
}
\caption{WF-10: Smart Scan Invoice Workflow.}
\label{fig:wf10}
\end{figure}

\section{Key Code Snippets with Explanation}

\subsection{Stockout Calculation Logic (n8n Code Node)}

The following JavaScript implements the stockout prediction formula within the WF-02 Code node:

\begin{lstlisting}[language=javascript, caption={WF-02 Stockout Calculation Code Node}, label={lst:stockout}]
const items = $input.all();
const today = new Date().toISOString().split('T')[0];

return items.map(item => {
  const { current_stock, avg_daily_sales } = item.json;
  const stock  = parseFloat(current_stock)  || 0;
  const sales  = parseFloat(avg_daily_sales) || 0;

  let days;
  if (stock === 0)    { days = 0; }
  else if (sales === 0) { days = 999; }
  else                { days = Math.floor(stock / sales); }

  const stockoutDate = new Date();
  stockoutDate.setDate(stockoutDate.getDate() + days);

  return {
    json: {
      ...item.json,
      days_to_stockout: days,
      estimated_stockout_date:
        days === 999 ? null : stockoutDate.toISOString().split('T')[0],
    }
  };
});
\end{lstlisting}

The \texttt{Math.floor} operation matches Equation~\ref{eq:stockout} precisely. The ternary
guard for \texttt{avg\_daily\_sales = 0} prevents division-by-zero errors in production.

\subsection{Supplier WSM Scoring Logic (n8n Code Node)}

\begin{lstlisting}[language=javascript, caption={WF-06 Weighted Sum Model Scoring}, label={lst:wsm}]
const suppliers = $input.all().map(i => i.json);

// Min-max normalise delivery time (faster = higher score)
const times = suppliers.map(s => s.delivery_time_days);
const minT  = Math.min(...times);
const maxT  = Math.max(...times);

return suppliers.map(s => {
  const R = parseFloat(s.reliability_score);
  const P = parseFloat(s.price_score);
  const D = maxT === minT ? 1 :
    (maxT - s.delivery_time_days) / (maxT - minT);

  const score = (R * 0.40 + P * 0.35 + D * 0.25) * 100;
  const grade = score >= 80 ? 'A' : score >= 65 ? 'B' :
                score >= 50 ? 'C' : 'D';

  return { json: { ...s, composite_score: score.toFixed(2),
                         supplier_grade: grade } };
});
\end{lstlisting}

\subsection{Supabase Realtime Subscription Pattern (React)}

All pages that display live data use the following pattern to subscribe to database changes:

\begin{lstlisting}[language=javascript, caption={Supabase Realtime Subscription Pattern}, label={lst:realtime}]
useEffect(() => {
  fetchProducts();          // initial data load
  const channel = supabase
    .channel('products-live')
    .on('postgres_changes',
        { event: '*', schema: 'public', table: 'Products' },
        () => fetchProducts())   // refetch on any change
    .subscribe();
  return () => supabase.removeChannel(channel);
}, []);
\end{lstlisting}

The cleanup function in the \texttt{useEffect} return ensures that the WebSocket channel is
properly removed when the component unmounts, preventing memory leaks.

\subsection{Inventory Health Score Calculation (Frontend)}

The dashboard calculates an overall inventory health score using the following formula:

\begin{equation}
\label{eq:health_score}
H = \min\!\left(100,\, \max\!\left(0,\, \left\lfloor \frac{\sum (n_i \cdot w_i)}{N} \right\rfloor\right)\right)
\end{equation}

where the numerator is $n_L \cdot 100 + n_M \cdot 70 + n_H \cdot 30 + n_O \cdot 5$; $n_L$, $n_M$, $n_H$, and $n_O$ are counts of LOW, MEDIUM, HIGH risk, and out-of-stock products respectively; and $N$ is the total number of products. This mathematical model is implemented in the frontend layer as shown in Listing~\ref{lst:health}.

\begin{lstlisting}[language=javascript, caption={Inventory Health Score Implementation}, label={lst:health}]
function computeHealthScore(products) {
  const N  = products.length;
  if (N === 0) return 100;
  const nL = products.filter(p => p.risk_level === 'LOW').length;
  const nM = products.filter(p => p.risk_level === 'MEDIUM').length;
  const nH = products.filter(p => p.risk_level === 'HIGH').length;
  const nO = products.filter(p => p.current_stock === 0).length;
  const weighted = nL*100 + nM*70 + nH*30 + nO*5;
  return Math.min(100, Math.max(0, Math.floor(weighted / N)));
}
\end{lstlisting}
